Here is one construction for each target. A bar denotes repeating digits, so $.\overline4=\dfrac{4}{9}$. Every expression on the right uses exactly four fours; factorials and roots introduce no additional numerals. The displayed target on the left is not part of the expression.\begin{enumerate}\item $0=(44-44)$.\item $1=((44-44))!$.\item $2=(4-\sqrt{(4.4-.4)})$.\item $3=(4-(4)^{(4-4)})$.\item $4=(4+((4-4))^{4})$.\item $5=(4+(4)^{(4-4)})$.\item $6=(4.4+(4\cdot .4))$.\item $7=(4+(4-((4-4))!))$.\item $8=(4-(.4-4.4))$.\item $9=(4+(4+((4-4))!))$.\item $10=\dfrac{44}{4.4}$.\item $11=(4+\dfrac{(4+(4)!)}{4})$.\item $12=((4)!-(4+(4+4)))$.\item $13=((4)!-\dfrac{44}{4})$.\item $14=(4+(4+(4+\sqrt{4})))$.\item $15=(4+\dfrac{44}{4})$.\item $16=(4+(4+(4+4)))$.\item $17=\dfrac{((4)!+44)}{4}$.\item $18=(.4+(4\cdot 4.4))$.\item $19=((4)!-(4+((4-4))!))$.\item $20=(.4-(4.4-(4)!))$.\item $21=((4)!-(4-((4-4))!))$.\item $22=(\sqrt{4}+(44-(4)!))$.\item $23=((4)!-(4)^{(4-4)})$.\item $24=(4+(44-(4)!))$.\item $25=((4)!+(4)^{(4-4)})$.\item $26=(\sqrt{4}+((4.4-.4))!)$.\item $27=(4+((4)!-((4-4))!))$.\item $28=(4+((4.4-.4))!)$.\item $29=(4+((4)!+((4-4))!))$.\item $30=\dfrac{(4+(4+4))}{.4}$.\item $31=((4)!+\dfrac{(4+(4)!)}{4})$.\item $32=\dfrac{(4)^{4}}{(4+4)}$.\item $33=((4)!+\dfrac{(4-.4)}{.4})$.\item $34=(44-\dfrac{4}{.4})$.\item $35=((4)!+\dfrac{44}{4})$.\item $36=(44-(4+4))$.\item $37=(\dfrac{(.4+(4)!)}{.4}-(4)!)$.\item $38=(44-(4+\sqrt{4}))$.\item $39=(44-\dfrac{\sqrt{4}}{.4})$.\item $40=((4)!+((4)!-(4+4)))$.\item $41=\dfrac{(.4+(4\cdot 4))}{.4}$.\item $42=(\sqrt{4}-(4-44))$.\item $43=(44-((4-4))!)$.\item $44=(44.4-.4)$.\item $45=(44+((4-4))!)$.\item $46=(4-(\sqrt{4}-44))$.\item $47=((4)!+((4)!-((4-4))!))$.\item $48=(4\cdot (4+(4+4)))$.\item $49=((4)!+((4)!+((4-4))!))$.\item $50=(4+(\sqrt{4}+44))$.\item $51=\dfrac{(.4-(4-(4)!))}{.4}$.\item $52=(4+(4+44))$.\item $53=(44+\dfrac{4}{.\overline{4}})$.\item $54=(44+\dfrac{4}{.4})$.\item $55=\dfrac{44}{(.4+.4)}$.\item $56=(4+(4+((4)!+(4)!)))$.\item $57=(\dfrac{(.4+(4)!)}{.4}-4)$.\item $58=\dfrac{((4)^{4}-(4)!)}{4}$.\item $59=(\dfrac{(4)!}{.4}-((4-4))!)$.\item $60=(44+(4\cdot 4))$.\item $61=(((4-4))!+\dfrac{(4)!}{.4})$.\item $62=\dfrac{(.4+(.4+(4)!))}{.4}$.\item $63=\dfrac{((4)^{4}-4)}{4}$.\item $64=((4)!-(4-44))$.\item $65=\dfrac{(4+(4)^{4})}{4}$.\item $66=((4)!-(\sqrt{4}-44))$.\item $67=(4+\dfrac{(4+(4)!)}{.\overline{4}})$.\item $68=(4+((4+4))^{\sqrt{4}})$.\item $69=\dfrac{(4-(.4-(4)!))}{.4}$.\item $70=(\sqrt{4}+((4)!+44))$.\item $71=\dfrac{((4)!+4.4)}{.4}$.\item $72=(4+((4)!+44))$.\item $73=\dfrac{((4)!+((4)!+\sqrt{.\overline{4}}))}{\sqrt{.\overline{4}}}$.\item $74=(4+\dfrac{(4+(4)!)}{.4})$.\item $75=\dfrac{(4)!}{(.4\cdot (.4+.4))}$.\item $76=((4\cdot ((4)!-4))-4)$.\item $77=((\dfrac{4}{.\overline{4}})^{\sqrt{4}}-4)$.\item $78=((4\cdot ((4)!-4))-\sqrt{4})$.\item $79=((4)!+\dfrac{((4)!-\sqrt{4})}{.4})$.\item $80=(4\cdot (44-(4)!))$.\item $81=((4-((4-4))!))^{4}$.\item $82=(\sqrt{4}-(4\cdot (4-(4)!)))$.\item $83=((4)!-\dfrac{(.4-(4)!)}{.4})$.\item $84=(4-(4\cdot (4-(4)!)))$.\item $85=((4)!+\dfrac{(.4+(4)!)}{.4})$.\item $86=(\dfrac{44}{.4}-(4)!)$.\item $87=((4)!+\dfrac{(4+(4)!)}{.\overline{4}})$.\item $88=(44+44)$.\item $89=((4)!+\dfrac{(\sqrt{4}+(4)!)}{.4})$.\item $90=((4\cdot (4)!)-(4+\sqrt{4}))$.\item $91=((4\cdot (4)!)-\dfrac{\sqrt{4}}{.4})$.\item $92=((4)!+((4)!+44))$.\item $93=((4\cdot (4)!)-\sqrt{\dfrac{4}{.\overline{4}}})$.\item $94=((4\cdot ((4)!-.4))-.4)$.\item $95=((4\cdot (4)!)-((4-4))!)$.\item $96=(4\cdot ((4.4-.4))!)$.\item $97=(((4-4))!+(4\cdot (4)!))$.\item $98=(.4+(4\cdot (.4+(4)!)))$.\item $99=\dfrac{44}{.\overline{44}}$.\item $100=\dfrac{44}{.44}$.\end{enumerate}Each equality follows by ordinary arithmetic; factorials are taken only of nonnegative integers. The repeating decimals are geometric series, which reduces their verification to rational arithmetic.
